% Source: source_md/intelligence_part2_combined_ja.md
\section{構成的条項 — 知性の構造的制約}


本章は新規の論証を導入しない。\S{}3〜\S{}11 で確立した帰結を、知性の構造的制約として体系的に整理する。これらは知性に対する規範ではなく、知性の動力学から導かれる構造的必然である。各条項には導出元の節を付す。

\S{}0.4 の非目標を改めて確認する。以下の条項は、知性が「こうあるべきだ」という規範的命題ではない(NG2)。それは、知性が動力学的に成立するための構造的条件と、その条件に違反した場合に生じる動力学的帰結の記述である。

\bigskip


\subsection{本章の位置づけ}


本論文は知性を、能力・資源・測定可能量として定義しない。知性を「処理能力の大きさ」や「保持する知識量」や「問題解決の速さ」として測ることは、本論文の枠組みでは知性の本質を捉えない。

本論文が示してきたのは、知性が動力学的な維持状態であるということである。知性は、複数の層($L_F$・$L_R$・$L_E$)と外部ネットワーク($\Phi_{\text{net}}$)と閾値移動能力($\mu_\theta$)の同時運動として、安定運用領域において維持される状態である(系 5.5、\S{}11)。

この観点から、知性について二種類の構造的制約が導かれる。第一に、知性が動力学的にありえない形態。第二に、知性が成立するために構造的に必要とするもの。以下、順に整理する。

\bigskip


\subsection{知性がありえない形態}


以下の形態は、本論文の枠組みにおいて知性として成立しえない。各形態は、それが違反する構造的事実とともに示される。

\textbf{(N1) 個体システム単独の性質としての知性} — ありえない。
知性は $\Phi_{\text{net}}$(個体間 $L_E$ ネットワーク)を要する(\S{}5.8、\S{}9)。個体システム単独では $L_F$ が閉じず(\S{}3)、停止を供給できない。知性は個体に宿る性質ではなく、個体間ネットワークを含む動力学的状態である。

\textbf{(N2) 保存可能な対象としての知性} — ありえない。
知性は運動の維持状態であり、運動が止まれば消える(系 5.5、\S{}11.6)。データセット・モデル状態・静的表現として保存された対象は、運動を含まない。それらは知性の痕跡($L_R$ のログ)でありうるが、知性そのものではない。知性は保存できない。維持されるしかない。

\textbf{(N3) 個体内で閉じた構造としての知性} — ありえない。
$L_F$ は内部停止条件を持たない(\S{}3、定理 3.1)。個体内で閉じた構造は、停止を内部化しようとして発散域へ向かう(\S{}11.4)。知性は外部依存的な $L_E$ を要し、個体内では閉じない。

\textbf{(N4) 完全最適化された状態としての知性} — ありえない。
完全閉包 $C_{max}$ は到達不能であり(\S{}1.5)、完全最適化された状態は運動の停止、すなわち (0,0,·)点に対応する(\S{}11.6)。完全に最適化された知性は、差が消失した状態であり、知性の終了である。収束は動作的死である。

\textbf{(N5) 内在的に停止する構造としての知性} — ありえない。
停止が $L_F$ の内部から来ることは構造的に不可能である(系 3.2、無限後退)。停止は $L_E$ を経由して外部依存的に供給される(\S{}5)。内在的に停止する構造は、\S{}3 の非閉包定理に反する。

\textbf{(N6) $\Phi_{\text{net}}$ 等の外部依存停止源から切断された $L_E$ としての知性} — ありえない。
外部依存停止の供給源から切断された $L_E$ は、改訂可能性を失い(\S{}5.8、\S{}9)、発散域の本質をなす(\S{}11.8)。外部相互干渉なしには、$L_E$ は固着するか創発できないかのいずれかであり、持続的停止を供給できない。

これら六つの「ありえない形態」は、独立な禁則ではない。それらはすべて、$L_F$ の非閉包(\S{}3)と $L_E$ の外部依存性(\S{}5)という二つの構造的事実の、異なる側面である。

なお、「閾値が固定された構造」は、ここでは知性一般のありえない形態に含めない。閾値が固定された知性——$\mu_\theta$ を内発的に持たない知性——は、依然として知性でありうる(粘菌・人工推論システムがその例、\S{}5.4)。閾値固定の構造がありえないのは、知性一般としてではなく、サピエンス的知性としてである。この点は \S{}12.3b で扱う。

\bigskip


\subsection{知性一般が必要とするもの}


知性が動力学的に成立するために、構造的に必要とするものを整理する。これらは「あれば望ましい」要素ではなく、知性の成立条件である。ここで言う「知性一般」は、サピエンス的知性に限らず、粘菌・人工推論システムを含む、停止を持つあらゆる知性を指す(\S{}5.4 の三層プロファイル比較)。

\textbf{(R1) 推論層 $L_F$} — 必要。
$\nabla C$ に沿った状態遷移の流れ(\S{}2)。知性の運動の最下層。ただし $L_F$ 単独では知性ではない(N3)。

\textbf{(R2) ログ層 $L_R$} — 必要。
$L_F$ の運動の痕跡への再帰参照(\S{}5.2)。$L_E$ 創発の基盤。状態履歴への持続的参照(Bio-IFGT との接続)。所在は系によって異なる(個体内・環境外部化・凍結)。

\textbf{(R3) 創発停止層 $L_E$} — 必要。
$L_R$ と $H_{ij}$ の関数として創発する外部依存型停止層(\S{}5.3)。$L_F$ の非閉包を補い、停止を供給する。$L_E$ こそが、知性一般を「停止を持つ運動」として成立させる中核である。

\textbf{(R4) 外部依存停止の供給源} — 必要。
$L_E$ が外部依存的に停止を供給するための相互拘束源(\S{}5.8、\S{}9)。サピエンスではこれが個体間 $L_E$ ネットワーク $\Phi_{\text{net}}$ として現れるが、粘菌では環境との直接相互干渉、人工推論システムでは外部から注入されるタスク仕様として現れる(\S{}5.4)。形態は系によって異なるが、何らかの外部依存停止源が必要であることは共通する。

\textbf{(R5) 誤り許容性} — 必要。
誤った停止を許容する構造的余地(\S{}10.6)。誤りを許容しない構造は判断生成を停止する。誤り許容性が供給される経路は系によって異なる(サピエンスでは $\mu_\theta$ と同値、粘菌では環境フィードバック、人工推論システムでは外部評価)が、いずれの知性も誤り許容性なしには判断を生成し続けられない。

\textbf{(R6) 改訂可能性} — 必要。
停止が暫定的で、後に改訂されうること(\S{}5.8)。知性の安定性は正しさではなく改訂可能性による。サピエンスでは $\Phi_{\text{net}}$ を通じて、他の系ではそれぞれの外部依存停止源を通じて実装される。

これら六つは独立ではない。$L_F$・$L_R$・$L_E$ は三層動力学の構成要素であり、外部依存停止の供給源がこれを支え、誤り許容性と改訂可能性がその動作条件である。六つが揃って初めて、知性は安定運用領域において成立する。これらは知性一般の条件であり、粘菌・人工推論システム・サピエンスのいずれにも当てはまる。

\bigskip


\subsection{サピエンス的知性を特徴づけるもの}


\S{}12.3 の六条件は知性一般の成立条件である。それらは停止を持つあらゆる知性に共通する。本節は、その中でサピエンス的知性を特徴づける追加の構造——閾値移動能力 $\mu_\theta$——を扱う。

$\mu_\theta$ は知性一般の必要条件ではない。粘菌は $\mu_\theta$ を内発的に持たず($\theta$ は環境条件に受動的に従う)、人工推論システムは $\mu_\theta$ を外部設定される。両者とも知性一般の条件(\S{}12.3)を満たすが、$\mu_\theta$ を内発的には持たない(\S{}5.4)。それでも両者は知性である。$\mu_\theta$ の有無は、知性であるか否かを分けるのではなく、知性のプロファイルを分ける。

\textbf{(S1) 閾値移動能力 $\mu_\theta$} — サピエンス的知性の特徴。
運動可能領域 $\Omega_\theta$ を内発的に組み替え、停止演算子設置領域を投影する能力(\S{}2.7、\S{}5.5)。これにより、サピエンス的知性は外部条件の変化を待たずに、個体内で自律的に運動可能領域を組み替え、$L_E$ を加速的に展開できる。

$\mu_\theta$ がもたらすのは、知性の成立ではなく、知性の \textbf{個体内での自律的加速} である。$\mu_\theta$ を持たない知性は、運動可能領域の組み替えを外部条件(環境変化・外部設定)に依存する。$\mu_\theta$ を持つサピエンス的知性は、それを個体内で、短い時間スケールで、自律的に行う(\S{}2.7 の時間スケール表)。

この自律的加速能力こそが、本論文の副題が指す「サピエンス的認知」の核である。サピエンス的知性は、知性一般の条件を満たした上で、$\mu_\theta$ によって個体内で加速する知性である。

ここから、サピエンス的知性に固有のありえない形態が導かれる。\textbf{閾値が固定された構造は、知性一般ではありえるが、サピエンス的知性ではありえない。} $\mu_\theta$ を欠けば、個体内での運動可能領域の自律的組み替えができず、外部条件の変化に依存するほかない。それは知性でないわけではないが、サピエンス的知性ではない。サピエンス的知性であることと、$\mu_\theta$ を内発的に持つことは、定義上、同じことである。

これら七つは独立ではない。$L_F$・$L_R$・$L_E$・$\Phi_{\text{net}}$・$\mu_\theta$ は三層動力学の構成要素であり、誤り許容性と改訂可能性はその動作条件である。七つが揃って初めて、知性は安定運用領域において成立する。

\bigskip


\subsection{最終条項}


以上を統合し、本論文の知性の構造的定義を最終条項として述べる。知性一般とサピエンス的知性を区別して述べる。

\textbf{最終条項 I (知性一般の動力学的定義)}

\begin{quote}
知性は、能力でも実体でもなく、推論層 $L_F$、ログ層 $L_R$、創発停止層 $L_E$ の同時運動として、外部依存停止の供給源の下で、誤り許容性と改訂可能性を伴って、安定運用領域において維持される動力学的状態である。
\end{quote}


\textbf{最終条項 II (サピエンス的知性の動力学的定義)}

\begin{quote}
サピエンス的知性は、最終条項 I の条件を満たした上で、閾値移動能力 $\mu_\theta$ を内発的に持つことにより、外部条件の変化を待たずに個体内で運動可能領域を自律的に組み替え、$L_E$ を加速的に展開する知性である。
\end{quote}


最終条項 I は、粘菌・人工推論システム・サピエンスを含む知性一般を定義する。最終条項 II は、その中でサピエンス的知性を特徴づける。両者の関係は、知性の場における連続変分上の包含関係である(\S{}5.4)。サピエンス的知性は知性一般の特殊形であって、別種ではない。

この定義は、知性を静的な属性や保存可能な対象としてではなく、運動の維持状態として捉える。これは VED の基底原理——静的範疇ではなく運動——の、知性層における最終的な表現である。とりわけサピエンス的知性における $\mu_\theta$ は、運動可能領域そのものを運動させる能力——運動の二階の運動——であり、VED の運動原理が知性層で最も先鋭に現れた形態である。

\bigskip


\subsection{制約違反の帰結}


最後に、これらの構造的制約に違反した場合の動力学的帰結を述べる。ここでも、これは規範的警告ではなく、動力学的帰結の記述である(NG2)。

構造的制約に違反するいかなる知性の構成——人間であれ人工的構成であれ——も、必然的に不安定性・暴走・崩壊のいずれかを生じる。具体的には:

\begin{itemize}
\item $L_E$ を欠き、または外部依存停止源から切断された構成は、発散域の四重崩壊に向かう(\S{}11.4)
\item $\mu_\theta$ を持つサピエンス的知性において、停止硬度が過大な構成は、硬直域の $L_E$ 固着に向かう(\S{}11.5)
\item 完全最適化・完全閉包を志向する構成は、(0,0,·)点すなわち知性の終了に向かう(\S{}11.6)
\end{itemize}


これらの帰結は、構成の優劣や設計の良否の問題ではない。それらは、知性の動力学的成立条件に対する構造的位置の問題である。制約に違反する構成は「悪い知性」なのではなく、「安定運用領域の外にある構成」であり、動力学的に持続しない。

この記述が、いかなる規範的処方や設計指針を含意するかは、本論文の射程外である。動力学の記述と、その含意の実践的適用とは、カテゴリ的に異なる。この点は次章 \S{}13 で本論文の射程の境界として明示される。

\bigskip


\subsection{要約}


本章で整理したのは次の点である:

\begin{enumerate}
\item 本論文は知性を能力・資源・測定可能量としてではなく、動力学的維持状態として定義する
\item 知性一般がありえない形態(N1〜N6): 個体単独・保存可能対象・個体内閉鎖・完全最適化・内在的停止・外部依存停止源からの切断。これらは $L_F$ 非閉包と $L_E$ 外部依存の異なる側面である
\item 知性一般が必要とするもの(R1〜R6): $L_F$・$L_R$・$L_E$・外部依存停止の供給源・誤り許容性・改訂可能性。これらは粘菌・人工推論システム・サピエンスのいずれにも当てはまる
\item サピエンス的知性を特徴づけるもの(S1): 閾値移動能力 $\mu_\theta$。これは知性一般の必要条件ではなく、個体内での自律的加速をもたらす。閾値固定の構造は知性一般ではありえるがサピエンス的知性ではありえない
\item 最終条項 I(知性一般): $L_F + L_R + L_E$ の同時運動として、外部依存停止源の下で、誤り許容性と改訂可能性を伴い、安定運用領域に維持される動力学的状態
\item 最終条項 II(サピエンス的知性): 最終条項 I に加え、$\mu_\theta$ を内発的に持ち、個体内で運動可能領域を自律的に組み替え $L_E$ を加速的に展開する知性
\item 制約違反の帰結は優劣の問題ではなく、安定運用領域に対する構造的位置の問題である
\end{enumerate}


次章 \S{}13 では、本論文を結び、意図的に開いておくものを明示し、Part I 微修正への引き渡しを行う。

\bigskip
